\documentclass[a4paper, fleqn]{cas-dc}
\renewcommand{\figurename}{Fig.}
\usepackage{makecell}
\usepackage{multirow}
\usepackage[utf8]{inputenc}
\usepackage[T1]{fontenc}
\usepackage[ruled, linesnumbered]{algorithm2e}
\usepackage{amsmath, amssymb, amsfonts}
\usepackage{graphicx}
\usepackage{booktabs}
\usepackage{array}
\usepackage{tabularx}
\usepackage{float}
\usepackage{xcolor}
\usepackage{pifont}% For checkmarks and crosses
\usepackage{xspace}% Required for \xspace command
\hypersetup{hypertexnames=false}
% If the frontmatter runs over more than one page
% use the longmktitle option.
%\documentclass[a4paper,fleqn,longmktitle]{cas-dc}

\usepackage[numbers]{natbib}
%\usepackage[authoryear]{natbib}
%\usepackage[authoryear,longnamesfirst]{natbib}

%%%Author macros
\providecommand{\csdef}[1]{\expandafter\def\csname #1\endcsname}
\def\tsc#1{\csdef{#1}{\textsc{\lowercase{#1}}\xspace}}
\tsc{WGM}
\tsc{QE}
%%%

% Uncomment and use as if needed
%\newtheorem{theorem}{Theorem}
%\newtheorem{lemma}[theorem]{Lemma}
%\newdefinition{rmk}{Remark}
%\newproof{pf}{Proof}
%\newproof{pot}{Proof of Theorem \ref{thm}}

\begin{document}
\let\WriteBookmarks\relax
\def\floatpagepagefraction{1}
\def\textpagefraction{.001}

% Short title
\shorttitle{}

% Short author
\shortauthors{Yin et al.}  

% Main title of the paper
\title [mode = title]{R3-RAG: Reweight, Rerank, and Reflect for Evidence-Calibrated Scientific Multimodal Reasoning}  

% Title footnote mark
% eg: \tnotemark[1]
%\tnotemark[1] 

% Title footnote 1.
% eg: \tnotetext[1]{Title footnote text}
%\tnotetext[1]{} 

% First author
%
% Options: Use if required
% eg: \author[1,3]{Author Name}[type=editor,
%       style=chinese,
%       auid=000,
%       bioid=1,
%       prefix=Sir,
%       orcid=0000-0000-0000-0000,
%       facebook=<facebook id>,
%       twitter=<twitter id>,
%       linkedin=<linkedin id>,
%       gplus=<gplus id>]

\author[1]{Zihao Yin}
\ead{yinzihao@stu.xju.edu.cn}

\author[1]{Tao Wang}
\cormark[1]
\ead{twang@xju.edu.cn}

\author[1]{Yurong Qian}
\ead{qyr@xju.edu.cn}

\author[1]{Kai Wang}
\ead{wangkai@xju.edu.cn}

\author[1]{Ping Hu}
\ead{huping@xju.edu.cn}

\author[1]{Li Zhang}
\ead{Zhangl@xju.edu.cn}

\affiliation[1]{
  organization={Xinjiang University},
  city={Urumqi},
  postcode={830049},
  country={China}
}

\cortext[1]{Corresponding author}
% Footnote text
% For a title note without a number/mark
%\nonumnote{}

% Here goes the abstract
\begin{abstract}
Multimodal Retrieval-Augmented Generation (MM-RAG) grounds vision-language models in external evidence, yet mainstream linear frameworks may struggle to adapt retrieval and fusion to the varying visual and textual requirements of scientific questions.
Specifically, current systems may suffer from blind fusion due to static modality weights, cross-modal misalignment when rerankers fail to detect inconsistent evidence, and open-loop retrieval errors that propagate into generation.
To address these limitations, Reweight, Rerank, and Reflect-Augmented Generation (R3-RAG) is proposed as a closed-loop multimodal RAG framework that coordinates three modules for evidence-calibrated scientific reasoning.
To support evidence-calibrated scientific reasoning, it integrates intent-aware dynamic fusion to adaptively align modal focus, coupled with an uncertainty-aware cascade reranker to filter potentially inconsistent cross-modal evidence.
Furthermore, the system employs a diagnose-rewrite-route reflection loop to support iterative correction of retrieval failures.
Experiments show that R3-RAG achieves a competitive 91.51\% accuracy on the ScienceQA benchmark and outperforms existing baselines on the MMMU test set official Science discipline subset, reaching 48.21\% accuracy.
The results suggest that the framework improves reliability by reducing non-dominant modality noise and by revisiting low-confidence retrieval results before final answer generation.
\end{abstract}

% Use if graphical abstract is present
%\begin{graphicalabstract}
%\includegraphics{}
%\end{graphicalabstract}

% Research highlights

%\begin{highlights}
 %\item Intent-aware dynamic fusion balances visual and textual evidence.
 %\item Uncertainty-aware cascade reranking reduces cross-modal hallucinations
 %\item A diagnose-rewrite-route loop enables self-correction of retrieval failures.
%\end{highlights}



% Keywords
% Each keyword is separated by \sep
\begin{keywords}

Autonomous Agents \sep

Intent-Aware Fusion \sep

Uncertainty-Aware Reranking \sep

Cognitive Reflection Loop \sep

Scientific Reasoning \sep

\end{keywords}

\maketitle

% Main text
\section{Introduction}\label{sec:introduction}

The rapid evolution of Multimodal Large Language Models (MLLMs) has expanded the processing of heterogeneous visual and textual data~\cite{yin2024survey}.
However, their application in knowledge-intensive domains, such as complex scientific reasoning, still faces significant reliability constraints~\cite{alampara2025probing}.
Large-scale parametric models have achieved strong performance across a wide range of generative tasks. However, relying solely on their internal memory can constrain these models, leading to static knowledge boundaries and factually incorrect hallucinations~\cite{Lewis2020, Ji2023}.
The introduction of RAG has helped mitigate these constraints by grounding generation in externally retrieved evidence for knowledge-intensive tasks~\cite{Lewis2020, Guu2020, Gao2023}.
Benchmarks such as ScienceQA~\cite{Lu2022} highlight that scientific discovery and complex reasoning problems often exhibit high degrees of heterogeneity and multimodal interaction: abstract physical laws usually exist in textual form, while specific structural details remain hidden within diagrams.
Improving scientific reasoning is important because it requires intelligent systems to operate over heterogeneous evidence sources, including unstructured text, structured knowledge graphs, and visual data~\cite{Liu2025}.
Moreover, such reasoning involves more than information acquisition alone; it also benefits from reflection mechanisms~\cite{Shinn2023, Madaan2023} that revisit complex logical conflicts, such as cases where visual observations appear to contradict textual common sense.
However, existing multimodal RAG frameworks remain largely constrained by static linear pipelines, lacking the dynamic calibration and iterative feedback needed to reconcile cross-modal inconsistencies effectively.
These limitations motivate the development of an adaptive multimodal RAG system that integrates dynamic intent perception, cross-modal evidence calibration, and iterative cognitive reflection.

Despite significant progress in multimodal RAG, many existing MRAG frameworks still rely on rigid pipelines, which struggle to adapt retrieval and fusion to diverse queries and often lead to suboptimal modality fusion and insufficient calibration of evidence relevance and consistency~\cite{Abootorabi2025, Ling2025, Mei2025}.
Specifically, mainstream MRAG architectures are predominantly built upon a static embedding matching paradigm.
Representative frameworks, such as MuRAG~\cite{Chen2022} and Uni-RAG~\cite{Li2025}, leverage pre-trained dual-tower or unified encoders (for example, CLIP~\cite{Radford2021}) to map visual and textual data into a shared semantic space.
The core advantage of this paradigm lies in breaking the static parametric boundaries of models, enabling the utilization of external heterogeneous data to supplement the generation context.
However, this one-size-fits-all linear strategy operates on the assumption that retrieved contexts are always reliable, ignoring the highly dynamic dependency between visual and textual information in complex scientific problems.
Consequently, this rigid pipeline exposes three recurring limitations corresponding to the retrieval, reranking, and reasoning phases.

First, in terms of retrieval, current methods fall into the blind fusion trap.
As revealed by Liang et al.~\cite{Liang2022} in their study of the Modality Gap, existing works typically rely on predefined frozen weights or coarse-grained global alignment.
This strategy inherently ignores geometric biases within the embedding space~\cite{Ethayarajh2019, Gao2021}.
For instance, as illustrated in Fig.~\ref{fig:vs}, when handling high-visual-dependency tasks such as circuit identification, static systems fail to suppress noise from the textual modality and may retrieve irrelevant electromagnetic theories.
Conversely, for pure theoretical derivations, the forced intervention of visual features introduces destructive background noise.
\begin{figure*}[t] 
  \centering

  \includegraphics[width=1.0\linewidth]{Figure_1.png}
  
\caption{
Illustration of structural bottlenecks in static Multimodal RAG pipelines.
}
  \label{fig:vs}
\end{figure*}

Second, the subsequent reranking stage may suffer from cross-modal misalignment.
Traditional rankers follow unimodal text retrieval logic and cannot quantify the logical consistency between images and text.
More critically, directly introducing VLMs for scoring is not a panacea.
As highlighted by the POPE study~\cite{Li2023} on Object Hallucination, VLMs are often subject to statistical priors, assigning high confidence scores to non-existent objects.
Even advanced models like VisRAG~\cite{Yu2025} may remain vulnerable to such high-confidence errors when uncertainty is not explicitly calibrated.
This can allow semantically overlapping but factually irrelevant documents to crowd out useful evidence and degrade the reasoning context.

Third, at the reasoning level, existing pipelines are often limited by open-loop execution~\cite{Asai2024}.
Most systems adopt a linear execution mode (i.e., one-shot retrieval, one-shot generation).
Within this unidirectional workflow, any noise introduced during the initial retrieval phase goes unchecked and is inevitably amplified throughout the reasoning chain, culminating in cascading hallucination errors~\cite{Yan2024}.
For example, in drug-interaction queries, an incorrect chemical structure retrieval can lead to the derivation of completely fabricated metabolic pathways.
Such logically consistent but factually incorrect generation is unacceptable in high-stakes scientific fields.

To address the above limitations of static RAG pipelines, this paper introduces an adaptive framework with a closed-loop retrieval refinement process.
Firstly, to overcome the blind fusion trap, an intent-aware dynamic fusion mechanism is proposed.
The proposed system does not rely solely on a fixed linear search process; instead, it first decodes the implicit modal intent within a user query, for example, by distinguishing between visual recognition needs and theoretical derivation needs, to actively adjust the focus of information fusion.
Subsequently, to address cross-modal misalignment, an adaptive gated cascade reranking module is constructed between retrieval and generation.
Unlike traditional coarse ranking based solely on semantic similarity, this module uses the predictive uncertainty of multimodal large models as a confidence signal to estimate the consistency between visual evidence and textual claims, thereby reducing highly semantically relevant but factually inconsistent noise.
Finally, when the generated confidence remains below a threshold, the framework activates a feedback mechanism to diagnose the likely failure mode, whether it originates from missing entities or modal conflicts, and reroutes the query among heterogeneous knowledge sources.

The main contribution of this work is not to propose another independent retriever or generator, but to formulate multimodal RAG as an evidence-calibration process. Unlike static pipelines that directly append retrieved evidence to the generation context, R3-RAG calibrates evidence at three stages: query-time modality reweighting, candidate-level uncertainty-aware reranking, and feedback-based retrieval correction. This three-stage calibration is particularly suitable for scientific multimodal reasoning, where visual and textual dependency varies substantially across questions and incorrect evidence can easily propagate into final answers.

Accordingly, R3-RAG is positioned as a problem-driven evidence-calibration framework for scientific multimodal reasoning, rather than as a general-purpose RAG architecture or a standalone algorithmic advance in the core retrieval-generation mechanism of RAG.

In summary, the main contributions of this paper are as follows:

\begin{itemize}
    \item \textbf{Intent-Aware Evidence Reweighting for Closed-Domain Multimodal Retrieval:} We propose an intent-guided dynamic alpha mechanism that converts query-level visual/textual dependency into an adaptive fusion weight. Unlike static CLIP-style fusion or fixed image-text interpolation, this mechanism shifts the hybrid query representation toward the dominant modality, reducing non-dominant modality noise in scientific questions with different visual-textual requirements.

    \item \textbf{Uncertainty-Gated Cross-Modal Evidence Calibration:} We design a two-stage cascade reranker that combines efficient textual coarse filtering with VLM-based visual logic auditing. Rather than directly trusting VLM relevance scores, the proposed gate uses the dispersion of VLM scores as a confidence proxy to decide how strongly visual alignment should influence final evidence ranking, thereby reducing high-relevance but factually inconsistent evidence.

    \item \textbf{Diagnose-Rewrite-Route Reflection over Heterogeneous Evidence Sources:} We introduce a lightweight reflection loop that diagnoses retrieval failures and routes revised queries among vector, graph, and multimodal retrieval sources. Compared with text-only reflection mechanisms, this loop explicitly considers modality-specific failures, such as missing visual evidence, entity-relation gaps, and keyword-level mismatches, enabling iterative correction before final answer generation.
\end{itemize}

The remainder of this paper is organized as follows. Section~\ref{sec:relatedwork} reviews related work on multimodal RAG, cross-modal reranking, and agentic reflection mechanisms. Section~\ref{sec:methodology} presents the proposed R3-RAG framework, Section~\ref{sec:experiments} reports the experimental setup and empirical analyses, and Section~\ref{sec:conclusion} concludes the paper.

\section{Related Work}\label{sec:relatedwork}
This section situates the proposed framework within three lines of research that are closely related to evidence-calibrated multimodal reasoning.
We first review multimodal RAG frameworks, then discuss cross-modal reranking techniques for evidence filtering, and finally examine agentic reflection and reasoning loops that motivate our closed-loop design.

\subsection{Multimodal RAG Frameworks}

Multimodal RAG extends the capability of Large Language Models (LLMs) by incorporating heterogeneous data sources, such as text, images, and structured knowledge graphs, to ground generation in external evidence.
Early approaches like REVEAL~\cite{Hu2023} demonstrated the potential of enhancing Vision-Language Models via decoupled retrievers, indicating that external visual evidence can aid reasoning.
To bridge the gap between modalities, representative frameworks such as MuRAG~\cite{Chen2022} and Uni-RAG~\cite{Li2025} pioneered the use of dual-tower architectures or unified embedding spaces to align image-text pairs into a shared semantic space.
Furthermore, BGE-M3~\cite{Chen2024} proposed a versatile embedding model to handle multi-granularity retrieval within a unified framework.
For instance, RA-CM3~\cite{Yasunaga2023} proposed a retrieval-augmented multimodal language model that interleaves text and images during both retrieval and generation phases, allowing for richer contextual grounding.
Additionally, ColPali~\cite{Faysse2025} revolutionized document retrieval by directly processing PDF pages as images using multi-vector representations, effectively preserving complex visual layouts and figures.
Moreover, WebQA-style frameworks~\cite{Chang2022} have been developed to perform multi-hop reasoning across independent image and text corpora, aggregating cross-modal evidence to synthesize answers.

However, these mainstream methods primarily rely on static fusion strategies.
They typically use frozen predefined encoders such as CLIP~\cite{Radford2021}, which inherently lack the flexibility to adapt to dynamic query intents.
As noted in studies on the Modality Gap~\cite{Liang2022}, such one-size-fits-all approaches struggle to regulate the balance between visual micro-details and textual theories, often leading to suboptimal retrieval in complex scientific tasks like ScienceQA~\cite{Lu2022}.
To address this, our work introduces an intent-guided dynamic $\alpha$ mechanism to achieve on-demand alignment.

\subsection{Cross-Modal Reranking}

Ensuring the logical consistency between retrieved multimodal candidates and the user query is critical for reducing hallucinations.
In the reranking phase, RagVL~\cite{Chen2025} focused on the Multi-granularity Noisy Correspondence (MNC) phenomenon, utilizing text-based rankers to filter candidates.
More recently, VisRAG~\cite{Yu2025} attempted to simplify the workflow by employing a fully visual pipeline, feeding retrieved images directly into VLMs for relevance assessment.
To further optimize cross-modal alignment, contemporary systems frequently adopt advanced architectures as fine-grained evaluators.
For example, BLIP-2~\cite{Li2023BLIP2} based Image-Text Matching (ITM) utilizes a lightweight Querying Transformer to bridge the modality gap, computing precise alignment scores between query text and visual patches.
To address the language priors inherent in such models, the MASS~\cite{Chung2025} framework explicitly mitigates the reliance on linguistic statistical biases, ensuring the reranker genuinely focuses on visual content, thereby drastically improving the visual accuracy of image-text matching.
More recently, enhanced visual-language foundation models like ELIP~\cite{Zhan2025} have refined the reranking stage by specifically optimizing for failure cases where small but semantically crucial visual objects are ignored during coarse retrieval, improving cross-modal retrieval robustness.

Nevertheless, existing reranking paradigms still face calibration limitations.
Text-based rankers often overlook fine-grained visual cues, while directly deploying VLMs for scoring is prone to statistical biases.
As revealed by the POPE~\cite{Li2023} study, VLMs suffer from Object Hallucination, often assigning high confidence scores to non-existent objects due to prior probabilities.
This lack of uncertainty calibration makes it difficult to intercept semantically relevant but factually contradictory noise.
To address this, the proposed Uncertainty-Aware Cascade Reranking module utilizes score dispersion as a proxy to filter high-confidence errors.

\subsection{Agentic Reflection and Reasoning Loops}

To overcome the limitations of rigid, linear retrieval pipelines, Agentic reasoning has gained traction for its potential to introduce self-correction capabilities.
Self-RAG~\cite{Asai2024} introduced special critique tokens to evaluate retrieval quality on demand, allowing the model to decide when to retrieve or generate.
Similarly, CRAG~\cite{Yan2024} incorporated triggers for external web searches to rectify local retrieval failures when the internal knowledge base is insufficient.
Additionally, Reflexion~\cite{Shinn2023} proposed dynamic memory updates via linguistic feedback, enabling systems to revise prior reasoning traces.
Moreover, frameworks like ReAct~\cite{Yao2023} interleave reasoning traces with task-specific actions, allowing agents to iteratively query databases and dynamically adjust search strategies based on intermediate observations.
For complex multi-step problems, Tree of Thoughts (ToT)~\cite{Yao2023ToT} enables models to concurrently explore multiple reasoning paths, evaluating local solutions and safely backtracking upon encountering dead ends.
Building upon this, Graph of Thoughts (GoT)~\cite{Besta2024} models the reasoning space as a complex network, synthesizing evidence from multiple branching paths via heuristic searches.

Despite these advances, current reflection mechanisms are predominantly confined to the textual domain.
Most multimodal systems still operate in an open-loop manner, lacking the ability to handle modality-specific failures, such as missing visual diagrams or missing textual definitions, and to recursively route queries across heterogeneous sources, including vector databases and knowledge graphs.
To address this gap, our paper proposes a multimodal diagnose-rewrite-route closed loop, which extends reflection from the textual domain to multimodal RAG and enables the system to diagnose and revise retrieval failures.

\textbf{Naming clarification.} A recent Findings of EMNLP 2025 paper also uses the acronym R3-RAG for reasoning with textual retrieval via reinforcement learning~\cite{li2025r3rag}. In this paper, R3-RAG instead denotes Reweight, Rerank, and Reflect: an inference-time closed-domain multimodal RAG framework for evidence-calibrated scientific QA, without additional parameter training. Thus, despite the acronym overlap, the two works differ in task scope and technical design.

\section{Methodology}\label{sec:methodology}

In complex multimodal scientific reasoning, the proposed framework improves the generation process $P(y \mid Q, V)$ by first calibrating the external evidence set $\mathcal{E}$ and then grounding answers $y$ in the calibrated evidence.
To achieve this goal, it extends the traditional static linear retrieval pipeline with a closed-loop Reweight-Rerank-Reflect workflow (Fig.~\ref{fig:framework}).
Rather than only appending retrieved documents, this architecture sequentially extracts multimodal evidence, filters cross-modal inconsistencies, and revises low-confidence retrieval results before yielding the final answer.
Guided by this workflow, the methodology unfolds in four interconnected stages: Section~\ref{sec:intent_retrieval} introduces Intent-Aware Dynamic Multimodal Retrieval to adaptively decouple modal focus and construct a high-recall candidate pool, thereby mitigating the "blind fusion" bottleneck; Section~\ref{sec:reranking} details the Uncertainty-Aware Cascade Reranking module, which uses VLM score dispersion to reduce semantically relevant but factually contradictory evidence; Section~\ref{sec:reflection_loop} elaborates on the Agentic Reflection Reasoning Loop, which implements the Diagnose-Rewrite-Route protocol for iterative error correction; and finally, Section~\ref{sec:algorithm_workflow} formalizes the integrated Algorithm Workflow.
\begin{figure*}[t] 
  \centering

  \includegraphics[width=1.0\linewidth]{Figure_2.png}
\caption{
Overview of the proposed closed-loop multimodal RAG framework. The system retrieves evidence from heterogeneous local sources, including NanoVectorDB for dense retrieval, LightRAG for graph retrieval, and OpenCLIP for multimodal image-text retrieval. Retrieved candidates are aggregated, calibrated by an uncertainty-aware two-stage cascade reranker, optionally refined through a diagnose-rewrite-route reflection loop, and finally passed to Qwen2.5-VL for evidence-grounded answer generation.
}
  \label{fig:framework}
\end{figure*}
\subsection{Intent-Aware Dynamic Multimodal Retrieval}\label{sec:intent_retrieval}

In multimodal scientific reasoning, traditional RAG systems relying on static fusion weights, such as $\alpha = 0.5$, often suffer from disruptive noise from non-dominant modalities, such as irrelevant visual backgrounds in purely theoretical tasks.
To bridge this Modality Gap, an intent-guided dynamic $\alpha$ calibration mechanism (Fig.~\ref{fig:semantic_intent_parsing}) is proposed to geometrically reconstruct the query's latent representation at inference time.
Specifically, the system first parses the user query to estimate a visual-dependency score $p_v(Q) \in [0, 1]$, and maps this score into a continuous adaptive fusion weight $\alpha_{dyn}(Q)$.
Subsequently, it utilizes $\alpha_{dyn}$ to interpolate the independent textual $\mathbf{e}_{txt}$ and visual $\mathbf{e}_{vis}$ embeddings into a unified Hybrid Query Vector $\mathbf{e}_{hybrid}$.
Geometrically, this dynamic fusion shifts the hybrid vector toward the dominant modality's feature region, thereby reducing non-dominant noise and improving the signal-to-noise ratio of the retrieved Top-$k$ multimodal evidence.
\begin{figure*}[t] 
  \centering
  \includegraphics[width=1.0\linewidth]{Figure_3.png}
\caption{
Detailed architecture of the Intent-Aware Dynamic Multimodal Retrieval module. Given a query $Q$ and visual context $V$, the semantic intent parser estimates a visual-dependency score $p_v(Q) \in [0, 1]$ according to Eq.~\eqref{eq:visual_dependency}. OpenCLIP encoders produce normalized textual and visual embeddings, fused by the continuous adaptive weight $\alpha_{dyn}(Q)$ computed from Eq.~\eqref{eq:dynamic_alpha}. The hybrid query vector $\mathbf{e}_{hybrid}$ in Eq.~\eqref{eq:hybrid_query} is used for Top-$15$ multimodal image-text retrieval to produce candidate evidence. The three panels respectively illustrate intent parsing (Section~\ref{sec:semantic_intent_parsing}), adaptive alpha fusion (Section~\ref{sec:dynamic_weight_calibration}), and the geometric intuition of vector rotation in the embedding space (Section~\ref{sec:geometric_vector_reconstruction}).
}
  \label{fig:semantic_intent_parsing}
\end{figure*}
\subsubsection{Semantic Intent Parsing}\label{sec:semantic_intent_parsing}

First, the mapping of the user query $Q$ into a continuous visual-dependency score is formulated.
To achieve this, a suite of domain-specific semantic trigger sets is defined: a visual-sensitive set $\mathcal{T}_{vis}$ comprising tokens such as \textit{look at}, \textit{diagram}, and \textit{color} that signify visual dependency, and a theoretical-sensitive set $\mathcal{T}_{theo}$ comprising tokens such as \textit{definition}, \textit{principle}, and \textit{formula} that denote textual dependency (each set contains approximately 50 tokens; the full lists are provided in the released code repository).

The visual-dependency score $p_v(Q) \in [0, 1]$ quantifies the degree to which query $Q$ relies on visual information:
\begin{equation}
p_v(Q) =
\begin{cases}
1.0, & \text{if } \exists t \in \mathcal{T}_{vis}, t \subset Q \\
0.0, & \text{if } \exists t \in \mathcal{T}_{theo}, t \subset Q \\
0.5, & \text{otherwise}
\end{cases}
\label{eq:visual_dependency}
\end{equation}

where $p_v(Q) = 1.0$ indicates strong visual dependency, $p_v(Q) = 0.0$ indicates textual-theoretical reasoning, and $p_v(Q) = 0.5$ represents a neutral, balanced state. Note that Eq.~\eqref{eq:visual_dependency} instantiates $p_v(Q)$ with three anchor values corresponding to discrete intent categories for reproducible experiments. In principle, $p_v(Q)$ can be implemented as any calibrated continuous estimator (e.g., a lightweight classifier outputting soft probabilities), and the downstream fusion mechanism in Eq.~\eqref{eq:dynamic_alpha} supports the full continuum $p_v(Q) \in [0, 1]$ without modification.
\subsubsection{Dynamic Weight Calibration}\label{sec:dynamic_weight_calibration}

Rather than selecting $\alpha$ from a discrete set, we define $\alpha_{dyn}$ as a continuous bounded coefficient controlled by the visual-dependency score $p_v(Q) \in [0, 1]$. This enables smooth interpolation between textual-dominant and visual-dominant retrieval:
\begin{equation}
\alpha_{dyn}(Q) = \alpha_{\min} + (\alpha_{\max} - \alpha_{\min}) \cdot p_v(Q)
\label{eq:dynamic_alpha}
\end{equation}
where $\alpha_{\min} = 0.1$ and $\alpha_{\max} = 0.9$ define the feasible interval $\alpha_{dyn} \in [0.1, 0.9]$. When $p_v(Q) = 0.0$, $\alpha_{dyn} = 0.1$ favors text; when $p_v(Q) = 0.5$, $\alpha_{dyn} = 0.5$ yields balanced fusion; when $p_v(Q) = 1.0$, $\alpha_{dyn} = 0.9$ favors visual features. Intermediate values (e.g., $0.37, 0.62, 0.78$) are supported by the linear mapping without modification.

Eq.~\eqref{eq:dynamic_alpha} replaces the static fusion convention $\alpha \equiv 0.5$ with a query-conditioned continuum, where $\alpha_{dyn}$ governs the rotation direction of the composite query vector along the text-visual geodesic arc (Section~\ref{sec:geometric_vector_reconstruction}). The clamped bounds $[0.1, 0.9]$ prevent the information collapse of a single modality, ensuring that the system retains a minimal feature representation of the secondary modality even under extreme intent scenarios.
\subsubsection{Geometric Vector Reconstruction}\label{sec:geometric_vector_reconstruction}

The final hybrid query vector $\mathbf{e}_{hybrid}$ is derived as the weighted sum of the normalized textual and visual embeddings:
\begin{equation}
\mathbf{e}_{hybrid} = \alpha_{dyn} \cdot \frac{\text{Enc}_{vis}(V)}{\|\text{Enc}_{vis}(V)\|} + (1 - \alpha_{dyn}) \cdot \frac{\text{Enc}_{txt}(Q)}{\|\text{Enc}_{txt}(Q)\|}
\label{eq:hybrid_query}
\end{equation}

Here, the denominator term $\|\cdot\|$ denotes $L_2$ norm normalization, designed to project heterogeneous visual and textual features onto a unified unit hypersphere, thereby eliminating magnitude discrepancies. Meanwhile, $\alpha_{dyn}$ serves as a linear interpolation coefficient that governs the rotation direction of the composite vector $\mathbf{e}_{hybrid}$ along the text-visual geodesic arc.
Specifically, a larger weight impels the query vector to shift towards the visual feature cluster center, thereby intensifying retrieval sensitivity towards visual details.

\begin{itemize}
    \item When $p_v(Q) \to 1$ (i.e., $\alpha_{dyn} \to 0.9$), the query vector rotates towards the visual centroid, prioritizing the capture of evidence with similar visual characteristics during Top-$k$ nearest neighbor retrieval.
    \item Conversely, when $p_v(Q) \to 0$ (i.e., $\alpha_{dyn} \to 0.1$), the query vector retreats towards the textual centroid, effectively filtering out distractors that are visually similar but semantically irrelevant.
    \item When $p_v(Q) = 0.5$ (i.e., $\alpha_{dyn} = 0.5$), the query vector maintains a balanced position, evenly weighting both modalities.
\end{itemize}

Through this dynamic calibration, the system supports on-demand retrieval and improves the signal-to-noise ratio (SNR) of heterogeneous evidence recall.
\subsection{Uncertainty-Aware Adaptive Two-Stage Cascade Reranking}\label{sec:reranking}

Although the intent-guided retrieval strategy adjusts the modal distribution of the candidate set, the initial retrieval pool $\mathcal{E}_{init}$ may still retain noise characterized by high semantic relevance but fine-grained factual inconsistencies.
To balance computational efficiency and cross-modal alignment precision, a Coarse-to-Fine Two-Stage Cascade Reranking architecture is proposed, incorporating an adaptive gating mechanism grounded in predictive uncertainty, as illustrated in Fig.~\ref{fig:uncertainty_gated_ceranking}.
\begin{figure*}[t] 
  \centering

  \includegraphics[width=1.0\linewidth]{Figure_4.png}
\caption{
Pipeline of the Uncertainty-Aware Adaptive Two-Stage Cascade Reranking module. The input is the initial evidence pool obtained from heterogeneous retrieval. Stage I (Section~\ref{sec:stage_textual_filtering}) uses BGE-Reranker-M3 as a cross-encoder to compute textual relevance scores via Eq.~\eqref{eq:textual_score} and perform Top-$15$ $\to$ Top-$5$ coarse filtering. Stage II (Section~\ref{sec:stage_visual_alignment}) uses Qwen2.5-VL-32B with rubric-based prompting to evaluate visual-textual consistency via Eq.~\eqref{eq:visual_score}. The dispersion of VLM scores in Eq.~\eqref{eq:visual_dispersion} controls the adaptive gate $\lambda$ in Eq.~\eqref{eq:adaptive_gate}, and the final ranking score is computed by fusing normalized textual and visual scores according to Eq.~\eqref{eq:final_score}.
}
  \label{fig:uncertainty_gated_ceranking}
\end{figure*}
\subsubsection{Stage I: Efficient Textual Coarse Filtering}\label{sec:stage_textual_filtering}

To manage large-scale retrieval candidates and reduce computational overhead, Stage I deploys a lightweight textual reranker serving as a coarse-grained filter.
Given the query $Q$ and the initial retrieved document set $\mathcal{E}_{init} = \{d_1, d_2, \dots, d_M\}$, the semantic relevance score is computed using a Cross-Encoder architecture (specifically, BGE-Reranker-M3):
\begin{equation}
s_{txt}^{(i)} = \text{CrossEnc}(Q, \text{Content}(d_i))
\label{eq:textual_score}
\end{equation}

where $s_{txt}^{(i)}$ denotes the semantic relevance score assigned to the $i$-th candidate document, while $\text{CrossEnc}(\cdot)$ represents the scoring function of a Transformer-based Cross-Encoder model.
Unlike bi-encoder architectures, this function facilitates full self-attention interaction between the query $Q$ and the document content $\text{Content}(d_i)$.
This mechanism targets fine-grained semantic matching signals often lost in vector embedding spaces, thereby providing preliminary filtering to alleviate the computational burden on the subsequent VLM inference.
\subsubsection{Stage II: Fine-Grained Visual Logic Alignment}\label{sec:stage_visual_alignment}

To bridge the visual understanding gap inherent in the text-based coarse filtering of Stage I, the second stage deploys a Multimodal Large Language Model (specifically, Qwen2.5-VL-32B) acting as a fine-grained evaluator.
Departing from coarse binary classification, a newly devised Rubric-based Prompting Strategy instructs the VLM to quantify the logical consistency of image-text evidence on a continuous scale of $[0.0, 10.0]$.
For each candidate $d_i \in \mathcal{E}_{cand}$ and its associated image $I_i$ retained from the previous stage, the VLM computes a visual alignment score:
\begin{equation}
s_{vis}^{(i)} = \text{VLM}(Q, I_i, \text{Content}(d_i); \mathcal{P}_{rubric})
\label{eq:visual_score}
\end{equation}

where $s_{vis}^{(i)}$ denotes the fine-grained cross-modal consistency score for the $i$-th candidate, while $\mathcal{P}_{rubric}$ represents the rubric-based prompt constraints injected into the VLM inference process.
Unlike standard likelihood scoring, this formulation requires the model to evaluate the multimodal triplet formed by the query $Q$, candidate image $I_i$, and text content $\text{Content}(d_i)$ against rubric-defined logical criteria, such as spatial relations and causal inference, thereby yielding a semantically calibrated scalar output.
\subsubsection{Variance-Based Dynamic Gating Mechanism}\label{sec:dynamic_gating}
Directly fusing the textual score $s_{txt}$ and visual score $s_{vis}$ can be affected by VLM calibration bias: when the VLM hallucinates, it may output indiscriminately high scores for all documents with negligible differences, reducing ranking discriminability.

To address this, the dispersion of the VLM score distribution serves as a proxy for its discriminative certainty.
Specifically, the standard deviation $\sigma_{vis}$ of the visual scores for the current batch of candidates is calculated:
\begin{equation}
\sigma_{vis} = \sqrt{\frac{1}{K_c} \sum_{i=1}^{K_c} (s_{vis}^{(i)} - \mu_{vis})^2}
\label{eq:visual_dispersion}
\end{equation}

where $\sigma_{vis}$ quantifies the dispersion of visual alignment scores within the candidate pool of size $K_c$, while $\mu_{vis}$ denotes the batch mean.
Serving as a proxy for uncertainty, this statistical moment evaluates the VLM's discriminative confidence: a low $\sigma_{vis}$ signals model indecision or mode collapse (failing to distinguish between candidates), whereas a high $\sigma_{vis}$ indicates that the model has captured fine-grained quality differences among the evidence, thereby supporting an increased visual weight in the subsequent gated fusion.

Based on $\sigma_{vis}$, an adaptive gating function $\lambda(\cdot)$ is constructed to dynamically regulate the visual weight:
\begin{equation}
\lambda = \text{Clamp} (\lambda_{base} + \gamma \cdot (\sigma_{vis} - \tau),\; \lambda_{min},\; \lambda_{max})
\label{eq:adaptive_gate}
\end{equation}

where $\lambda$ denotes the adaptive gating weight employed for final score fusion, derived by modulating the base weight $\lambda_{base}$ with the deviation of visual dispersion $\sigma_{vis}$ from a preset threshold $\tau$.
Here, $\gamma$ serves as a sensitivity coefficient governing the gain of the gating mechanism in response to uncertainty fluctuations.
In our experiments, we set $\lambda_{base} = 0.3$, $\gamma = 2.0$, $\tau = 0.5$, and clamp $\lambda$ to $[\lambda_{min}, \lambda_{max}] = [0.1, 0.9]$.
When $\sigma_{vis} > \tau$, the visual contribution increases; when $\sigma_{vis} < \tau$, indicating that the VLM is indecisive, the gate falls back toward textual retrieval.
The $\text{Clamp}$ function ensures the fusion ratio remains within numerically stable bounds to prevent retrieval collapse under extreme conditions.

The final fusion score $S_{final}^{(i)}$ is calculated as follows:
\begin{equation}
S_{final}^{(i)} = (1 - \lambda) \cdot \text{Norm}(s_{txt}^{(i)}) + \lambda \cdot \text{Norm}(s_{vis}^{(i)})
\label{eq:final_score}
\end{equation}

where $S_{final}^{(i)}$ denotes the final multimodal fusion evidence score for the $i$-th candidate, while $\text{Norm}(\cdot)$ represents the normalization operation applied to eliminate scale discrepancies between heterogeneous scoring spaces.
This formulation provides dynamic interpolation between textual and visual signals via the adaptive gate $\lambda$, prioritizing visual logic when the VLM exhibits high discriminative power (larger $\lambda$) and relying more on textual semantic retrieval when the VLM's visual discriminative confidence is low, thereby improving ranking stability under multimodal conflicts.
\subsection{Agentic Reflection Reasoning Loop}\label{sec:reflection_loop}
Traditional RAG systems often follow a Retrieve-then-Generate unidirectional linear paradigm.
Consequently, once the initial retrieval is inaccurate, whether due to local matching errors or cross-modal semantic drift, the system has limited opportunity to revise the retrieved context.
To address this limitation, an Agentic Reflection Reasoning Loop is proposed, as illustrated in Fig.~\ref{fig:agentic_reflection_loop}.
This mechanism reformulates the retrieval pipeline as an iterative state refinement process. Governed by a Diagnose-Rewrite-Route protocol, the system updates both query representations and retrieval strategies when confidence is low.
\begin{figure*}[t]
  \centering
  \includegraphics[width=1.0\linewidth]{Figure_5.png}
\caption{
Workflow of the Agentic Reflection Reasoning Loop. At each iteration, the system evaluates the calibrated confidence of the reranked evidence using Eq.~\eqref{eq:final_score}. When $\max(S_{final}) < \tau_{reflect}$ or conflicting evidence is detected according to the trigger mechanism in Section~\ref{sec:reflection_trigger}, the diagnostic agent identifies the retrieval failure type $d_{err}$ in Eq.~\eqref{eq:error_descriptor}, such as ``Visual Evidence Missing'', ``Entity Relational Gap'', or ``Keyword Mismatch''. Conditioned on $d_{err}$, the query $Q_t$ is rewritten as $Q_{t+1}$ via Eq.~\eqref{eq:query_rewrite} and routed to the corresponding retrieval source: $\mathcal{K}_{vec}$ (vector database), $\mathcal{K}_{graph}$ (knowledge graph), or $\mathcal{K}_{mm}$ (multimodal pool). The loop terminates when the confidence exceeds the threshold or the maximum number of iterations $T_{max}$ is reached, after which the calibrated evidence is passed to final generation.
}
  \label{fig:agentic_reflection_loop}
\end{figure*}
\subsubsection{Reflection Trigger Mechanism}\label{sec:reflection_trigger}
At each time step $t$, the system maintains a reasoning state $\mathcal{S}_t = (Q_t, \mathcal{E}_t, c_t)$, where $Q_t$ represents the current query, $\mathcal{E}_t$ denotes the retrieved evidence set, and $c_t$ indicates the current confidence level.

The calibrated confidence score $S_{final}$ derived in Section~\ref{sec:reranking} is leveraged as the system's intrinsic self-evaluation metric.
The reflection loop is activated upon the satisfaction of either of the following criteria:

\begin{itemize}
    \item \textbf{Low Confidence Alert:} $\max(S_{final}) < \tau_{reflect}$, where $\tau_{reflect} = 0.5$ after score normalization, indicating that the currently retrieved evidence may be insufficient to support reliable reasoning.
    \item \textbf{High Information Entropy:} The evidence set $\mathcal{E}_t$ exhibits logical dissonance, for example, when the Top-$1$ and Top-$2$ candidates present contradictory factual statements, making convergence to a unified conclusion less reliable.
\end{itemize}
\subsubsection{Diagnosis-Rewrite-Route Protocol}\label{sec:diagnose_rewrite_route}

When reflection mode is activated, a lightweight diagnostic agent $\mathcal{A}_{meta}$ is used to execute three sequential atomic operations:

\noindent\textbf{1. Diagnosis.}

The meta-agent analyzes the current failure state to generate a specific error descriptor $d_{err}$. This process goes beyond simple binary classification, offering a fine-grained diagnosis grounded in a predefined error taxonomy:
\begin{equation}
\begin{aligned}
    d_{err} &= \mathcal{A}_{meta}(\text{Diagnose} \mid Q_t, \mathcal{E}_t), \\
    d_{err} &\in \mathcal{D}_{err}, \\
    \mathcal{D}_{err} &=
    \left\{
    \begin{array}{l}
    \text{``Visual Evidence Missing''},\\
    \text{``Entity Relational Gap''},\\
    \text{``Keyword Mismatch''}
    \end{array}
    \right\}.
\end{aligned}
\label{eq:error_descriptor}
\end{equation}

Each error type directly determines the subsequent routing action: ``Visual Evidence Missing'' triggers multimodal retrieval to recover overlooked diagrams; ``Entity Relational Gap'' targets the knowledge graph for multi-hop structural reasoning; and ``Keyword Mismatch'' falls back to dense vector retrieval with expanded queries.

\noindent\textbf{2. Rewrite.}
Conditioned on the diagnostic output $d_{err}$, the system employs a generator to map the original query $Q_t$ into a refined semantic space, yielding a revised query $Q_{t+1}$:
\begin{equation}
Q_{t+1} = \mathcal{M}_{gen}(\text{Rewrite} \mid Q_t, d_{err})
\label{eq:query_rewrite}
\end{equation}

This process facilitates a directed semantic shift of the query, aligning it more closely with the latent distribution of the knowledge base.
For example, it effectively transforms an ambiguous natural language question into a structured query containing specific entity relations.

\noindent\textbf{3. Route.}
This step constitutes the key routing decision in the closed loop.
The system dynamically selects the subsequent retrieval action $a_{t+1}$ based on the error type, within the defined action space $\mathcal{A} = \{ \mathcal{K}_{vec}, \mathcal{K}_{graph}, \mathcal{K}_{mm} \}$:
\begin{itemize}
    \item \textbf{Route A (Visual-Intensive Route):} If $d_{err}$ is ``Visual Evidence Missing'', the system routes the query to the Multimodal Pool ($\mathcal{K}_{mm}$). In this scenario, the system not only retrieves image-text pairs but also recomputes the visual-dependency score with an elevated weight ($p_v(Q) \to 1.0$), so that visual features dominate the retrieval process.
    
    \item \textbf{Route B (Structural-Logic Route):} If $d_{err}$ is ``Entity Relational Gap'', the system targets the Knowledge Graph ($\mathcal{K}_{graph}$), activating multi-hop subgraph retrieval to complete the missing logical reasoning chain.

    \item \textbf{Route C (Semantic-Broad Route):} If $d_{err}$ is ``Keyword Mismatch'', the system falls back to the Vector Database ($\mathcal{K}_{vec}$), utilizing an expanded query for broader semantic search.
\end{itemize}
\subsection{Algorithm Workflow}\label{sec:algorithm_workflow}

To enhance the reproducibility and logical clarity of the proposed framework, Algorithm~\ref{alg:R3_RAG} encapsulates the methodologies detailed in Sections 3.1 through 3.3 into algorithmic pseudocode.
\begin{algorithm}[t]

    \caption{R3-RAG Inference Pipeline} 
    \label{alg:R3_RAG}
    
   
    \SetKwInOut{Input}{Input}
    \SetKwInOut{Output}{Output}

    \Input{Query $Q$, Image $I$, Knowledge Base $\mathcal{K}$}
    \Output{Answer $\hat{y}$}

    
    \textbf{Initialize:} Evidence set $\mathcal{E} \leftarrow \emptyset$, Retrieval source $\mathcal{K}_{target} \leftarrow \mathcal{K}$, Iteration $t \leftarrow 0$, Max iterations $T_{max} \leftarrow 3$\;


    \While{$t < T_{max}$}{
        \tcp{Step 1: Intent-Aware Retrieval (Sec 3.1)}
        $p_v \leftarrow \text{CalcVisualDependency}(Q)$\;
        $\alpha_{dyn} \leftarrow 0.1 + 0.8 \cdot p_v$\;
        $\mathcal{E}_{init} \leftarrow \text{Retrieve}(Q, \mathcal{K}_{target}, \alpha_{dyn})$\;

        \tcp{Step 2: Adaptive Gated Reranking (Sec 3.2)}
        $\mathcal{E}_{ranked}, \sigma_{vis} \leftarrow \text{Rerank}(Q, \mathcal{E}_{init})$\;
        $\lambda \leftarrow \text{CalcGate}(\sigma_{vis})$\;
        $score_{conf} \leftarrow \text{WeightedScore}(\mathcal{E}_{ranked}, \lambda)$\;

        \tcp{Step 3: Reflection Check (Sec 3.3)}
        \eIf{$score_{conf} \ge \tau_{reflect}$}{
            $\mathcal{E} \leftarrow \mathcal{E}_{ranked}$\;
            \textbf{break}\;
        }{
            $d_{err} \leftarrow \text{Diagnose}(Q, \mathcal{E}_{ranked})$\;
            $\mathcal{E} \leftarrow \mathcal{E}_{ranked}$\;
            $Q \leftarrow \text{Rewrite}(Q, d_{err})$\;
            $\mathcal{K}_{target} \leftarrow \text{Route}(d_{err})$\;
            $t \leftarrow t + 1$\;
        }
    }
    
   
    $\hat{y} \leftarrow \text{Generate}(Q, I, \mathcal{E})$\;
    \Return{$\hat{y}$}
\end{algorithm}
\section{Experiments}\label{sec:experiments}

This section evaluates whether the proposed R3-RAG framework improves scientific multimodal reasoning under both in-domain and more challenging out-of-distribution test settings. Section~\ref{sec:experimental_setup} describes the experimental setup, and Section~\ref{sec:main_results} presents the main comparative results on ScienceQA~\cite{Lu2022} and MMMU~\cite{Yue2024MMMU}. Sections~\ref{sec:statistical_reliability}--\ref{sec:ablation_study} then report statistical reliability, computational cost, qualitative case studies, and ablation experiments.
We evaluate R3-RAG on ScienceQA~\cite{Lu2022} and the MMMU~\cite{Yue2024MMMU} test set official Science discipline subset, using a unified large-model-driven multimodal pipeline in which large models support logic routing, visual perception, multimodal reranking, and final answer summarization via Ollama. Key settings include an 8,192-token context window, adaptive fusion weights ($\alpha \in [0.1, 0.9]$), and distinct decoding temperatures for deterministic routing ($T_{route} = 0.1$) versus generative voting ($T_{vote} = 0.7$). All experiments were conducted on a single node equipped with a single NVIDIA A100 GPU.
\subsection{Experimental Setup}\label{sec:experimental_setup}
\textbf{Datasets}. We rigorously evaluate R3-RAG on two multimodal benchmarks to assess both in-domain performance and out-of-distribution (OOD) generalization.
Our primary evaluation is conducted on the ScienceQA~\cite{Lu2022} benchmark.
This dataset comprises 21,208 multimodal multiple-choice questions covering natural, social, and language sciences across Grades 1--12.
We report accuracy on the official test split, which contains 4,241 samples.
The dataset is characterized by its high heterogeneity, with approximately 48.7\% of questions containing image contexts and 48.2\% of questions containing text contexts.
This diversity serves as an ideal testbed for evaluating the system's capability in K-12 multimodal scientific reasoning.

To further probe the system's behavior on college-level tasks, we incorporate the official Science discipline subset from the MMMU~\cite{Yue2024MMMU} test set.
This subset is categorized by subject and covers Agriculture, Biology, Chemistry, Geography, Math, and Physics, yielding 2,713 examples in total.
Unlike the K-12 level ScienceQA, MMMU~\cite{Yue2024MMMU} presents higher-order reasoning challenges requiring more specialized domain knowledge.
We emphasize that this subject-level official Science discipline subset is used as a diagnostic out-of-distribution evaluation, rather than as a replacement for the full-MMMU benchmark. Therefore, the resulting scores should not be directly compared with full-MMMU leaderboard results unless the same subset and evaluation protocol are used. To ensure strict comparability and avoid unfair comparisons with full-MMMU leaderboard results, the zero-shot MMMU baselines in Table~\ref{tab:mmmu_results} were re-evaluated locally on this exact subset, while the few-shot Qwen2.5-VL result is reported as an additional prompting reference.

\textbf{Implementation Details}. All experiments were conducted on a single NVIDIA A100 GPU. We employed a unified large-model-driven multimodal pipeline, where large models supported intent decomposition, routing, visual perception, multimodal reranking, and final answer summarization.
The hybrid retrieval module integrates three sources: NanoVectorDB for dense retrieval, LightRAG for graph traversal, and OpenCLIP for multimodal retrieval. Key hyperparameters include an 8,192-token context window, an adaptive fusion weight $\alpha \in [0.1, 0.9]$, and a coarse-to-fine reranking strategy (Top-$15$ $\to$ Top-$5$).
Decoding temperatures are set to $T_{route} = 0.1$ for deterministic routing and $T_{vote} = 0.7$ for generative voting (samples $n = 3$).

\textbf{Retrieval corpus and leakage control}. The local retrieval corpus is constructed before evaluation from filtered non-evaluation scientific background materials, including textual passages, curated image-text pairs, and entity-relation triples for graph retrieval.
We keep the corpus disjoint from the official ScienceQA test split and the MMMU test-set Science subset, and do not index evaluation questions, answer labels, answer keys, lecture fields, explanations/rationales, or near-duplicate question texts from the evaluated splits.
No external web retrieval or online evidence is used during evaluation.

Accordingly, RagVL, Self-RAG, and CRAG are evaluated under the same closed-domain Qwen2.5-VL-32B protocol as R3-RAG, using the same local evidence sources and excluding external web retrieval.
Other LLM/VLM/tool rows in Table~\ref{tab:main_results} are retained as representative reported results from the ScienceQA leaderboard and prior literature.
\begin{table*}[t]
    \centering
    \caption{Top-$1$ answer accuracy comparison on the ScienceQA dataset (Accuracy \%). \#P denotes the reported model/backbone parameter scale when available. Categories include: NAT (Natural Science), SOC (Social Science), LAN (Language Science), TXT (Text Context), IMG (Image Context), NO (No Context), G1--6 (Grades 1--6), and G7--12 (Grades 7--12). The R3-RAG result in this main comparison table corresponds to the seed-21 run; the five-seed mean and confidence interval are reported separately in Table~\ref{tab:seed_reliability}. The RagVL, Self-RAG, CRAG, and R3-RAG rows are locally implemented under a closed-domain Qwen2.5-VL-32B setting; other rows are representative reported results from the ScienceQA leaderboard and related literature. Web-augmented methods using external online evidence, such as HM-RAG, are discussed separately in Section~\ref{sec:scienceqa_results} because they correspond to a different retrieval-access setting from the closed-domain comparison considered in this table.}
    \label{tab:main_results}
    \vspace{1.5mm} 

    \resizebox{\textwidth}{!}{
        \begin{tabular}{c|l|c|ccc|ccc|cc|c}
            \toprule
            \multirow{2}{*}{\textbf{Learning}} & \multirow{2}{*}{\textbf{Models}} & \multirow{2}{*}{\textbf{\#P}} & \multicolumn{3}{c|}{\textbf{Subject}} & \multicolumn{3}{c|}{\textbf{Context Modality}} & \multicolumn{2}{c|}{\textbf{Grade}} & \multirow{2}{*}{\textbf{Average}} \\
            \cline{4-11}
             &  &  & \textbf{NAT} & \textbf{SOC} & \textbf{LAN} & \textbf{TXT} & \textbf{IMG} & \textbf{NO} & \textbf{G1--6} & \textbf{G7--12} &  \\
            \midrule
            \multicolumn{1}{c|}{Baseline} & Human & - & 90.23 & 84.97 & 87.48 & 89.60 & 87.50 & 88.10 & 91.59 & 82.42 & 88.40 \\
            \midrule
            \multirow{6}{*}{Zero-shot LLMs} & CoT GPT-4~\cite{Lu2023Chameleon} & 1T+ & 85.48 & 72.44 & 90.27 & 82.65 & 71.49 & 92.89 & 86.66 & 79.04 & 83.99 \\
             & HoT-T5\_Large~\cite{Yao2023HoT} & 738M & 84.46 & 79.08 & 84.64 & 82.89 & 75.81 & 88.15 & 83.88 & 82.47 & 83.38 \\
             & DDCoT (GPT-3)~\cite{Zheng2023} & 175B & 78.60 & 73.90 & 80.45 & 77.27 & 69.96 & 82.93 & 80.65 & 73.50 & 78.09 \\
             & CoT GPT-3 + Doc~\cite{Hsieh2023} & 173B & - & - & - & - & - & - & - & - & 79.91 \\
             & DDCoT (ChatGPT)~\cite{Zheng2023} & 175B & 80.15 & 76.72 & 82.82 & 78.89 & 72.53 & 85.02 & 82.86 & 75.21 & 80.15 \\
             & GPT-3 (0-shot)~\cite{Lu2022} & 173B & 75.04 & 66.59 & 78.00 & 74.24 & 65.74 & 79.58 & 76.36 & 69.87 & 74.04 \\
            \midrule
            \multirow{8}{*}{Zero-shot VLMs} & LaVIN-13B~\cite{Luo2023LaVIN} & - & - & - & - & - & - & - & - & - & 77.54 \\
             & LLaMA-SciTune~\cite{Horawalavithana2024} & 7B & 84.50 & \textbf{94.15} & 82.91 & 88.35 & 83.64 & 88.74 & 85.05 & 85.60 & 86.11 \\
             & LG-VQA (BLIP-2)~\cite{Ghosal2023} & - & - & - & - & - & - & - & - & - & 86.32 \\
             & Qwen2.5-VL~\cite{Bai2025Qwen25VL} & 7B & - & - & - & - & - & - & - & - & 67.10 \\
             & DeepSeek-VL~\cite{Lu2024DeepSeekVL} & 7B & 78.51 & 90.21 & 70.55 & 87.07 & 81.36 & 70.94 & 82.60 & 72.25 & 78.90 \\
             & LLaVA-NeXT~\cite{Liu2024LLaVANeXT} & 13B & 74.02 & 69.97 & 70.18 & 81.88 & 68.57 & 71.92 & 76.14 & 65.06 & 72.18 \\
             & LG-VQA (CLIP)~\cite{Ghosal2023} & - & - & - & - & - & - & - & - & - & 87.22 \\
             & Chat-UniVi~\cite{Jin2024} & 7B & 58.61 & 61.08 & 61.82 & 57.33 & 58.25 & 61.39 & 62.04 & 56.23 & 59.96 \\
            \midrule
            \multirow{3}{*}{\shortstack{Tool-LLM}} & 	Chameleon (ChatGPT)~\cite{Lu2023Chameleon} & 175B+ & 81.62 & 70.64 & 84.00 & 79.77 & 70.80 & 86.62 & 81.86 & 76.53 & 79.93 \\
             & CoT GPT-3 + Doc~\cite{Hsieh2023} & 173B & - & - & - & - & - & - & - & - & 79.91 \\
             & Chameleon (GPT-4)~\cite{Lu2023Chameleon} & 1T+ & 89.83 & 74.13 & 89.82 & 88.27 & 77.64 & 92.13 & 88.03 & 83.72 & 86.54 \\
             \midrule
            \multirow{3}{*}{\shortstack{Closed-domain\\local RAG}} & RagVL~\cite{Chen2025} & 32B & 86.06 & 83.01 & 85.00 & \textbf{95.94} & 79.33 & 87.39 & 87.33 & 81.21 & 85.15 \\
             & Self-RAG~\cite{Asai2024} & 32B & 85.52 & 89.65 & 87.91 & 87.58 & 84.48 & 90.24 & 90.93 & 79.96 & 87.01 \\
             & CRAG~\cite{Yan2024} & 32B & 71.00 & 73.79 & 89.09 & 68.62 & 66.04 & 89.06 & 78.45 & 72.38 & 76.28 \\
            \midrule
            \shortstack{Closed-domain\\multi-agent RAG} & \multirow{2}{*}{\textbf{R3-RAG}} & \multirow{2}{*}{32B} & \multirow{2}{*}{\textbf{91.34}} & \multirow{2}{*}{92.35} & \multirow{2}{*}{\textbf{91.18}} & \multirow{2}{*}{89.86} & \multirow{2}{*}{\textbf{90.98}} & \multirow{2}{*}{\textbf{93.17}} & \multirow{2}{*}{\textbf{93.25}} & \multirow{2}{*}{\textbf{88.40}} & \multirow{2}{*}{\textbf{91.51}} \\
             & & & & & & & & & & & \\
            \bottomrule
        \end{tabular}}
\end{table*} 

\begin{table*}[t]
    \centering

    \caption{Performance comparison (\%) on the MMMU~\cite{Yue2024MMMU} test set official Science discipline subset, grouped by subject. The overall score is computed over 2,713 examples across Agriculture, Biology, Chemistry, Geography, Math, and Physics. Zero-shot baselines were re-evaluated locally on this specific subset, and the few-shot Qwen2.5-VL result is reported as an additional prompting reference. These scores are not intended to be directly comparable with full-MMMU leaderboard results.}
    \label{tab:mmmu_results}
    \vspace{2mm} 
    \renewcommand{\arraystretch}{1.25} 
    
    \resizebox{\textwidth}{!}{
        \begin{tabular}{l|cccccc|c}
            \toprule 
            \textbf{Method} & \textbf{Agriculture} & \textbf{Biology} & \textbf{Chemistry} & \textbf{Geography} & \textbf{Math} & \textbf{Physics} & \textbf{Overall} \\
            \midrule 
            Qwen2.5-VL-32B~\cite{Bai2025Qwen25VL} & 47.04 & 53.04 & 39.14 & 46.90 & 44.36 & 47.06 & 45.52 \\
            Few-shot Qwen2.5-VL-32B & \textbf{51.22} & \textbf{53.91} & \textbf{40.80} & 48.32 & 43.37 & 48.77 & 46.81 \\
            RagVL~\cite{Chen2025} & 44.25 & 50.14 & 40.30 & 41.24 & 40.20 & 45.10 & 42.87 \\
            Self-RAG~\cite{Asai2024} & 45.64 & 49.28 & 34.99 & 48.14 & 37.03 & 44.36 & 42.46 \\
            CRAG~\cite{Yan2024} & 49.13 & 48.41 & 28.36 & 36.99 & 18.81 & 29.17 & 33.25 \\
            \textbf{R3-RAG (Ours)} & 47.74 & 44.93 & 38.64 & \textbf{49.73} & \textbf{50.10} & \textbf{61.03} & \textbf{48.21} \\
            \bottomrule 
        \end{tabular}
    }
\end{table*}
\subsection{Main Results}\label{sec:main_results}
In this section, we evaluate R3-RAG against representative baselines reported in prior work and against locally implemented closed-domain RAG baselines.
We first analyze the comparative results on the ScienceQA benchmark (Table~\ref{tab:main_results}), focusing on performance differences across subjects, modalities, and grade levels.
Subsequently, we evaluate the method on the MMMU~\cite{Yue2024MMMU} test set official Science discipline subset (Table~\ref{tab:mmmu_results}) to examine preliminary generalization behavior in higher-difficulty scientific questions. Beyond reporting overall accuracy, this section also discusses the comparative advantages and limitations of R3-RAG relative to closely related multimodal and self-reflective RAG schemes.

\subsubsection{Results on ScienceQA Dataset}\label{sec:scienceqa_results}
Table~\ref{tab:main_results} compares R3-RAG with representative ScienceQA baselines, including reported zero-shot LLM/VLM methods, reported tool-augmented LLMs, and locally implemented closed-domain RAG methods. As shown in Table~\ref{tab:main_results}, R3-RAG achieves an average accuracy of 91.51\%. Under the same closed-domain Qwen2.5-VL-32B protocol, R3-RAG exceeds Self-RAG (87.01\%) by 4.50 percentage points. Compared with the locally implemented RagVL and CRAG baselines, R3-RAG improves the average accuracy by 6.36 and 15.23 percentage points, respectively. These results indicate that the proposed reweight-rerank-reflect pipeline improves overall robustness under the evaluated closed-domain setting.

To provide a more fine-grained analysis, Figs.~\ref{fig:subject_performance}--\ref{fig:grade_level} compare R3-RAG with selected representative baselines, including Human, CoT GPT-4, LLaMA-SciTune (7B), and Chameleon (GPT-4), from three perspectives: subject category, context modality, and grade level. As shown in Fig.~\ref{fig:subject_performance}, R3-RAG achieves 91.34\% on NAT, 92.35\% on SOC, and 91.18\% on LAN, showing balanced performance across disciplines in this selected comparison group. Fig.~\ref{fig:context_modality} shows the modality-level comparison, where R3-RAG obtains 90.98\% on IMG, 89.86\% on TXT, and 93.17\% on NO. Although RagVL reports a higher TXT score (95.94\%) than R3-RAG (89.86\%), R3-RAG remains more balanced across heterogeneous context settings. As illustrated in Fig.~\ref{fig:grade_level}, R3-RAG achieves 93.25\% on G1--6 and 88.40\% on G7--12, showing consistent performance across grade levels.

Compared with closely related RAG schemes under the same local closed-domain backbone, the advantage of R3-RAG lies not only in its final accuracy but also in its evidence-calibration process before generation. RagVL mainly strengthens multimodal retrieval and reranking, but it does not explicitly revise low-confidence retrieval failures through a closed-loop mechanism. Self-RAG focuses on retrieval self-assessment, and CRAG performs corrective retrieval. In contrast, R3-RAG coordinates intent-aware modality reweighting, uncertainty-gated cross-modal reranking, and diagnose-rewrite-route reflection over heterogeneous local evidence sources. Through this coordinated pipeline, the framework suppresses non-dominant modality noise, filters semantically relevant but cross-modally inconsistent evidence, and revises low-confidence retrieval results before final answer generation.

We also note that HM-RAG~\cite{Liu2025} reports 93.73\% accuracy on ScienceQA using web retrieval and online evidence, whereas R3-RAG is evaluated under a closed-domain setting with vector, graph, and multimodal sources. These two settings serve different deployment needs: web retrieval provides broader coverage, while closed-domain retrieval offers controllable, reproducible evidence and does not depend on network availability---properties that matter for offline scientific environments, privacy-sensitive applications where external web access may be restricted, and scenarios where evidence provenance and corpus scope need to be auditable. Within the closed-domain setting, R3-RAG achieves a competitive 91.51\%, indicating that the evidence-calibration pipeline is effective without external search.
\begin{figure*}[htbp]
    \centering

    \includegraphics[width=\textwidth]{Figure_6.png}
    \caption{Performance comparison with selected representative baselines across different ScienceQA subjects.}
    \label{fig:subject_performance}
\end{figure*}
\begin{figure*}[htbp] 
    \centering
    \includegraphics[width=\textwidth]{Figure_7.png}
    \caption{Performance comparison with selected representative baselines based on context modalities.}
    \label{fig:context_modality}
\end{figure*} 
\begin{figure*}[htbp] 
    \centering
    \includegraphics[width=\linewidth]{Figure_8.png}
    \caption{Performance comparison with selected representative baselines across different grade levels (G1--6 and G7--12).}
    \label{fig:grade_level}
\end{figure*} 

\subsubsection{Results on the MMMU Test Set Official Science Discipline Subset}\label{sec:mmmu_results_subsection}
To assess preliminary generalization on more specialized scientific tasks, we evaluate R3-RAG on the official Science discipline subset from the MMMU test set, grouped by subject, following the diagnostic protocol described in Section~\ref{sec:experimental_setup}. As shown in Table~\ref{tab:mmmu_results}, R3-RAG achieves an overall accuracy of 48.21\% (1308/2713), compared with 46.81\% for few-shot Qwen2.5-VL-32B, 45.52\% for Qwen2.5-VL-32B, 42.87\% for RagVL, 42.46\% for Self-RAG, and 33.25\% for CRAG under the same subject-level evaluation protocol. This corresponds to improvements of 1.40, 2.69, 5.34, 5.75, and 14.96 percentage points over few-shot Qwen2.5-VL-32B, Qwen2.5-VL-32B, RagVL, Self-RAG, and CRAG, respectively.

At the subject level, the results reveal a sharper pattern: R3-RAG outperforms the baselines on Geography (49.73\%), Math (50.10\%), and Physics (61.03\%)---subjects that frequently involve diagrams, spatial reasoning, or multi-step structured derivations where reconciling visual and textual evidence matters. However, on Agriculture (47.74\% vs.\ 51.22\%), Biology (44.93\% vs.\ 53.91\%), and Chemistry (38.64\% vs.\ 40.80\%), few-shot Qwen2.5-VL-32B achieves higher accuracy. Notably, on Biology, R3-RAG also underperforms the zero-shot Qwen2.5-VL-32B (44.93\% vs.\ 53.04\%, a gap of $-$8.11 points), indicating that the evidence-calibration pipeline can introduce noise rather than reduce it in certain domains.

We attribute this subject-level divergence to two factors. First, MMMU Biology and Chemistry questions in the official Science discipline subset often require recall of specialized terminology (e.g., biochemical pathway names, IUPAC nomenclature) that resembles closed-book factual retrieval. When the local knowledge base lacks precise coverage of such fine-grained domain facts, the retrieval stage may return topically adjacent but factually irrelevant candidates, and the subsequent reranking and reflection stages may fail to reject them when the VLM lacks the specialized knowledge to detect the mismatch. In such cases, the retrieval pipeline adds uninformative context that can distract the backbone, leading to lower accuracy than directly answering from parametric knowledge. Second, Physics and Math questions in MMMU more frequently pair textual problem statements with visual diagrams (force diagrams, circuit schematics, geometric figures), creating a modality gap that the intent-aware fusion and cascade reranker are designed to address. This contrast suggests a boundary condition: R3-RAG benefits questions where cross-modal evidence calibration adds discriminative signal, but may degrade performance when a question primarily demands closed-book factual recall and the retrieval corpus provides insufficient coverage, in which case the retrieval overhead introduces distracting context without compensating gain.

We treat this boundary characterization as a finding rather than a limitation to be concealed. It indicates that evidence-calibration frameworks should be deployed selectively based on the expected modality dependency and knowledge coverage of the target domain, and motivates future work on corpus-adaptive gating that can decide whether to retrieve at all. The relatively low absolute scores across all evaluated methods further confirm that the MMMU official Science discipline subset remains challenging for current multimodal RAG systems. The computational cost associated with the closed-loop design is analyzed separately in Section~\ref{sec:cost_efficiency_analysis}.

\subsection{Statistical Reliability and Group-wise Analysis}\label{sec:statistical_reliability}
To assess the statistical reliability of the reported ScienceQA results, we repeat the full evaluation with five random seeds: 13, 21, 42, 87, and 100.
The main ScienceQA comparison in Table~\ref{tab:main_results} reports the seed-21 run, while Table~\ref{tab:seed_reliability} summarizes performance across all five seeds. Across these runs, R3-RAG obtains a mean accuracy of 91.24\% with a standard deviation of 0.32, and the corresponding 95\% confidence interval is [90.85, 91.64].
This narrow interval indicates that the observed performance is stable under stochastic decoding and routing variations.
Since ScienceQA is a multiple-choice benchmark, we further report answer-option class-wise precision, recall, and F1 scores by aggregating predictions across the five seeds.
The class-wise results in Table~\ref{tab:seed_reliability} show that classes A--D remain close in F1 score, while class E has substantially lower support and should therefore be interpreted cautiously.
Table~\ref{tab:group_reliability} additionally reports group-wise accuracy across modality, grade, and subject categories, showing that the method maintains relatively stable performance across heterogeneous ScienceQA subsets.

\begin{table*}[t]
    \centering
    \caption{Statistical reliability over five random seeds and answer-option class-wise performance on ScienceQA. Class-wise support is aggregated over five runs.}
    \label{tab:seed_reliability}
    \vspace{1mm}
    \begin{minipage}{0.43\textwidth}
        \centering
        \begin{tabular*}{\linewidth}{@{\extracolsep{\fill}}lcc@{}}
            \toprule
            \textbf{Seed} & \textbf{Correct/Total} & \textbf{Accuracy (\%)} \\
            \midrule
            13 & 3886/4241 & 91.63 \\
            21 & 3881/4241 & 91.51 \\
            42 & 3854/4241 & 90.87 \\
            87 & 3866/4241 & 91.16 \\
            100 & 3861/4241 & 91.04 \\
            \midrule
            Mean $\pm$ Std & -- & 91.24 $\pm$ 0.32 \\
            95\% CI & -- & [90.85, 91.64] \\
            \bottomrule
        \end{tabular*}
    \end{minipage}
    \hfill
    \begin{minipage}{0.53\textwidth}
        \centering
        \begin{tabular*}{\linewidth}{@{\extracolsep{\fill}}lcccc@{}}
            \toprule
            \textbf{Class} & \textbf{Support} & \textbf{Precision} & \textbf{Recall} & \textbf{F1} \\
            \midrule
            A & 8500 & 89.29 & 92.99 & 91.10 \\
            B & 8445 & 92.77 & 89.82 & 91.27 \\
            C & 2915 & 91.84 & 90.39 & 91.11 \\
            D & 1305 & 92.22 & 88.97 & 90.56 \\
            E & 40 & 57.14 & 60.00 & 58.54 \\
            \bottomrule
        \end{tabular*}
    \end{minipage}
\end{table*}

\begin{table}[t]
    \centering
    \caption{Group-wise accuracy stability across five random seeds on ScienceQA.}
    \label{tab:group_reliability}
    \vspace{1mm}
        \begin{tabular*}{\linewidth}{@{\extracolsep{\fill}}lcc@{}}
            \toprule
            \textbf{Group} & \textbf{N/run} & \textbf{Accuracy $\pm$ Std} \\
            \midrule
            IMG & 2017 & 90.40 $\pm$ 0.77 \\
            TXT & 789 & 88.98 $\pm$ 0.58 \\
            NO & 1435 & 93.67 $\pm$ 0.25 \\
            Grades 1--6 & 2724 & 92.93 $\pm$ 0.33 \\
            Grades 7--12 & 1517 & 88.21 $\pm$ 0.61 \\
            Language science & 1100 & 92.09 $\pm$ 0.37 \\
            Natural science & 2252 & 90.96 $\pm$ 0.40 \\
            Social science & 889 & 90.89 $\pm$ 0.79 \\
            \bottomrule
        \end{tabular*}
\end{table}

\subsection{Computational Cost and Efficiency Analysis}\label{sec:cost_efficiency_analysis}
Since R3-RAG is an inference-time framework and does not require additional parameter training, this section focuses on inference-time computational cost. To address practical deployment concerns, we instrumented the inference process on the full ScienceQA evaluation split of 4,241 examples.
The profiling run records wall-clock time, per-example latency, retrieval/reflection calls, token consumption, and stage-level runtime. Wall-clock time denotes the elapsed time of the full profiling run and is used to compute throughput, whereas average and P95 latency summarize the per-sample latency logs; therefore, the two timing views are reported separately rather than being forced into a strictly serial equivalence.
For R3-RAG, the profiled run obtains 91.63\% accuracy, which is close to the main result in Table~\ref{tab:main_results}; therefore, the following analysis reflects the same operating regime while exposing detailed runtime behavior.
At inference time, the overall runtime can be approximated as $T_{\mathrm{R3}}=T_{\mathrm{dec}}+T_{\mathrm{ret}}+T_{\mathrm{rerank}}+\mathbb{I}_{\mathrm{ref}}T_{\mathrm{ref}}+n_{\mathrm{sc}}T_{\mathrm{gen}}$, where $T_{\mathrm{dec}}$, $T_{\mathrm{ret}}$, $T_{\mathrm{rerank}}$, $T_{\mathrm{ref}}$, and $T_{\mathrm{gen}}$ denote decomposition, retrieval, reranking, reflection, and final generation costs, respectively.
With fixed candidate budgets and self-consistency samples in our setting, the dominant practical factor is the number of expensive generation/reflection calls rather than the retrieval operation itself.

\begin{table*}[t]
    \centering
    \caption{Efficiency comparison on the full ScienceQA evaluation split under 32B settings. Wall-clock time is the elapsed profile time used to compute throughput, while Avg. Latency and P95 are computed from per-sample latency logs. A dash indicates that the corresponding metric was not logged by that baseline profiler.}
    \label{tab:efficiency_comparison}
    \vspace{1mm}
    \resizebox{\textwidth}{!}{
        \begin{tabular}{lccccccccl}
            \toprule
            \textbf{Method} & \textbf{Accuracy} & \textbf{Wall-clock} & \textbf{Avg. Latency} & \textbf{P95} & \textbf{Throughput} & \textbf{VLM Calls} & \textbf{Reflection Calls} & \textbf{Avg. Tokens} & \textbf{Dominant Cost} \\
             & \textbf{(\%)} & \textbf{(h)} & \textbf{(s/ex.)} & \textbf{(s)} & \textbf{(ex./h)} & \textbf{(/ex.)} & \textbf{(/ex.)} & \textbf{(/ex.)} &  \\
            \midrule
            CRAG~\cite{Yan2024} & 76.28 & 6.11 & 6.01 & 10.11 & 694.2 & 1.00 & 0.00 & 1072.06 & VLM answering \\
            Self-RAG~\cite{Asai2024} & 87.01 & 30.92 & 26.25 & 39.9 & 137.2 & 1.43 & 1.00 & 721.01 & Reflection (61.2\%) \\
            RagVL~\cite{Chen2025} & 85.15 & 5.84 & 4.96 & 7.74 & 725.9 & 1.00 & 0.00 & 448.25 & Generation (99.6\%) \\
            \textbf{R3-RAG (profiled)} & \textbf{91.63} & 36.66 & 31.12 & 40.43 & 115.68 & 2.23 & 0.468 & 1653.19 & Final answering (88.14\%) \\
            \bottomrule
        \end{tabular}}
\end{table*}

\begin{table}[t]
    \centering
    \caption{Stage-level latency breakdown for the profiled R3-RAG run on ScienceQA.}
    \label{tab:r3_cost_breakdown}
    \vspace{1mm}
    \resizebox{\linewidth}{!}{
        \begin{tabular}{lcc}
            \toprule
            \textbf{Stage} & \textbf{Sec./example} & \textbf{Share} \\
            \midrule
            Decompose & 0.38 & 1.23\% \\
            Retrieval & 0.43 & 1.38\% \\
            Rerank & 1.62 & 5.21\% \\
            Reflection & 1.24 & 3.99\% \\
            Summary/final answering & 27.42 & 88.14\% \\
            \bottomrule
        \end{tabular}}
\end{table}

\begin{figure*}[t]
    \centering
    \begin{minipage}[t]{0.49\textwidth}
        \centering
        \textbf{(a)}\\[-1mm]
        \makebox[\linewidth][c]{\includegraphics[height=0.25\textheight, keepaspectratio]{Figure_9.png}}
    \end{minipage}
    \hfill
    \begin{minipage}[t]{0.49\textwidth}
        \centering
        \textbf{(b)}\\[-1mm]
        \makebox[\linewidth][c]{\includegraphics[height=0.25\textheight, keepaspectratio]{Figure_10.png}}
    \end{minipage}
    \caption{Computational efficiency visualization on the ScienceQA evaluation split. (a) Accuracy-latency trade-off among profiled RAG baselines. (b) Stage-level runtime breakdown of R3-RAG, showing that final summarization dominates the inference cost.}
    \label{fig:cost_efficiency}
\end{figure*}

As shown in Table~\ref{tab:efficiency_comparison}, R3-RAG is more computationally expensive than the three profiled comparison methods, but it also achieves the highest accuracy among them in the same full-set profiling protocol.
Fig.~\ref{fig:cost_efficiency} visualizes the corresponding accuracy-latency trade-off and the internal runtime distribution.
The additional latency mainly comes from the final summarization and self-consistency stage rather than the retrieval stack.
Table~\ref{tab:r3_cost_breakdown} shows that retrieval accounts for only 0.43 seconds per example (1.38\%), while final answering accounts for 27.42 seconds per example (88.14\%).
Inside the retrieval stage, graph, vector, and multimodal retrieval consume 0.213, 0.128, and 0.090 seconds per example, respectively, indicating that hybrid retrieval contributes limited runtime overhead compared with generation-side inference.
Compared with Self-RAG, whose reflection stage occupies 16.07 seconds per example and 61.2\% of total latency, R3-RAG invokes reflection conditionally (0.468 reflection calls per example), reducing the reflection-stage share to 3.99\%.
The profiled R3-RAG run uses 48.02/49.97 GB mean/max GPU memory, 85.21\% mean GPU utilization, and an estimated 10.97 kWh of GPU energy.
These results suggest that the current system is suitable for high-accuracy offline or server-side scientific QA scenarios, while latency-sensitive deployment should prioritize reducing self-consistency samples and applying the final vision-summary call only under high uncertainty.

\textbf{Scalability discussion.} The current profiling analysis is conducted under a fixed candidate budget (Top-$15$ $\to$ Top-$5$) on a single NVIDIA A100 GPU. The stage-level breakdown shows that the dominant cost comes from final self-consistency generation, which accounts for 88.14\% of the per-example latency, while retrieval and reranking account for a much smaller portion of the runtime. This suggests that, under the current setting, scalability is primarily constrained by generation-side inference rather than local evidence retrieval. As the corpus size or candidate budget increases, the retrieval and reranking stages are expected to become more influential, but they can benefit from established acceleration techniques such as approximate nearest-neighbor search, cached embeddings, batched reranking, and parallel retrieval across vector, graph, and multimodal sources. A systematic study of scaling behavior with respect to corpus size, candidate budget, and hardware parallelism is left for future work.

\subsection{Qualitative Analysis}\label{sec:qualitative_analysis}

\begin{figure*}[t!] 
  \centering
  \includegraphics[width=0.85\linewidth, height=0.35\textheight, keepaspectratio]{Figure_11.png}
  \caption{
    \textbf{Case Analysis under Visual Retrieval Noise.} 
    For a map-based state identification question, the retrieval pipeline fetches visually irrelevant candidates, such as plants and birds. The \textit{Cascade Reranker} in Section~\ref{sec:reranking} assigns a low relevance score ($0.15 < \tau_{reflect}$), rejects the noisy context, and lets the model rely on intrinsic geographic knowledge to identify the highlighted state.
  }
  \label{fig:robustness_case_1}
\end{figure*}


\begin{figure*}[t!] 
  \centering
  \includegraphics[width=0.85\linewidth, height=0.35\textheight, keepaspectratio]{Figure_12.png}
  \caption{
    \textbf{Case Analysis under Domain-Specific Semantic Drift.} 
    For a Punnett-square reasoning question, the text-oriented retrieval path drifts toward irrelevant agricultural content, such as crop harvesting theory. The Cascade Reranker in Section~\ref{sec:reranking} assigns a low relevance score ($0.27 < \tau_{reflect}$), rejects the misleading evidence, and the model resolves the genetic probability through self-consistency inference.
  }
  \label{fig:robustness_case_2}
\end{figure*}

\begin{figure*}[t!] 
  \centering
  \includegraphics[width=0.85\linewidth, height=0.35\textheight, keepaspectratio]{Figure_13.png}
  \caption{
    \textbf{Case Analysis under Cross-Modal Misalignment.} 
    When the retrieval pipeline fetches misaligned textual candidates, such as irrelevant chunks matching material or path, the Cascade Reranker in Section~\ref{sec:reranking} assigns a low relevance score ($-0.54 < \tau_{reflect}$) and rejects these candidates. This gated filtering mechanism triggers the reflection loop in Section~\ref{sec:reflection_trigger}, leading the Adaptive Alpha Fusion in Section~\ref{sec:dynamic_weight_calibration} to prioritize the visual modality ($p_v(Q) \to 1.0$, $\alpha_{dyn} \to 0.9$) and rely on zero-shot visual reasoning, for example, by identifying a herringbone pattern via self-consistency inference.
  }
  \label{fig:robustness_case_3}
\end{figure*}
To illustrate the behavior of the proposed R3-RAG framework under retrieval failures, three representative scenarios are analyzed, covering visual hallucinations, domain-specific semantic drift, and cross-modal misalignment.

As illustrated in Fig.~\ref{fig:robustness_case_1}, a geography task is first examined to show the system's response to visual retrieval noise.
The task requires identifying a highlighted state on a U.S. map.
Initially, the Adaptive Alpha Fusion module correctly identifies the visual-dominant intent ($p_v(Q) \to 1.0$, $\alpha_{dyn} \to 0.9$), prioritizing map features over text.
However, due to corpus sparsity, the retrieval pipeline suffers semantic drift, fetching irrelevant visual hallucinations such as plants and birds.
The Cascade Reranker flags these candidates with low relevance scores ($0.15 < \tau_{reflect}$) and triggers the Agentic Reflection Loop.
Although iterative query rewriting fails to retrieve valid map data, the system's explicit rejection of the noisy context reduces the risk of evidence-induced hallucination.
This gated filtering mechanism encourages the model to disregard external interference and rely on its intrinsic parametric knowledge, identifying the state as Wyoming through self-consistency inference.

Fig.~\ref{fig:robustness_case_2} further analyzes a genetics case characterized by text-dominant intent and domain-specific semantic drift.
The task requires calculating the probability of a homozygous recessive trait based on a Punnett square.
In this instance, the Adaptive Alpha Fusion module prioritizes text keywords ($p_v(Q) \to 0.0$, $\alpha_{dyn} \to 0.1$) to retrieve biological definitions.
Nevertheless, the retrieval pipeline fetches irrelevant agricultural noise, such as crop harvesting theory.
The Cascade Reranker intercepts this semantic drift with a low score ($0.27 < \tau_{reflect}$), activating the reflection loop.
Following unsuccessful external retrieval attempts, the system explicitly rejects the misleading agricultural context.
Through self-consistency inference, despite a minor disagreement detected in one reasoning sample, the majority vote infers the probability to be 0/4 (Option C).

Finally, to illustrate the framework's behavior under cross-modal alignment failures, Fig.~\ref{fig:robustness_case_3} presents a vision-dominant material identification task.
The question asks for the material of a path shown in an image.
The Adaptive Alpha Fusion module detects a strong visual dependency ($p_v(Q) \to 1.0$, $\alpha_{dyn} \to 0.9$), suppressing unreliable textual noise.
During retrieval, generic textual cues such as ``material'' and ``path'' lead to a cross-modal alignment failure, yielding a negative relevance score (BGE = -0.54, well below the 0.8 threshold).
Consequently, this uninformative textual noise is explicitly rejected, triggering the Agentic Reflection Loop.
After diagnosing the semantic ambiguity of the vague terms, the system falls back to visual reasoning. By analyzing visual patterns, such as the staggered arrangement of the tiles and their rough textures, all three reasoning samples converge on the correct answer (Option B: brick).

Collectively, these cases illustrate how the system can block uninformative interference, whether visual, domain-specific, or textual, and rely on internal reasoning when external retrieval provides limited value.
\begin{figure*}[htbp] 
    \centering
    \includegraphics[width=0.72\linewidth, height=0.28\textheight, keepaspectratio]{Figure_14.png} 
    \caption{Waterfall chart illustrating the incremental accuracy improvements on the ScienceQA dataset.}
    \label{fig:ablation_waterfall}
\end{figure*} 
\subsection{Ablation Study}\label{sec:ablation_study}

\subsubsection{Effectiveness of Key Components}
To examine the contribution of our three proposed core modules, namely Adaptive Alpha Fusion, Cascade Reranking, and the Agentic Reflection Loop, we conducted an incremental ablation study on the ScienceQA dataset.
As illustrated in the waterfall chart (Fig.~\ref{fig:ablation_waterfall}) and detailed in Table~\ref{tab:ablation_full}, we decompose the performance gains step-by-step.
First, the baseline Vector-Graph RAG achieves an initial accuracy of 85.95\%.
Introducing the Adaptive Alpha Fusion improves overall performance to 87.33\% (a gain of +1.38\%), with notable improvements on visually-dependent categories: SOC gains 8.68 points (82.23\% $\to$ 90.91\%) and IMG gains 3.77 points (82.90\% $\to$ 86.67\%).
However, this gain is not uniform: Language Science (LAN) drops from 87.82\% to 84.55\% ($-$3.27), and No-Context (NO) drops from 90.03\% to 87.01\% ($-$3.02).
These degradations on text-only and language-heavy subsets are consistent with the boundary condition observed on MMMU Biology (Section~\ref{sec:mmmu_results_subsection}): when a question lacks visual context, forcing the query representation toward a visual modality---even minimally---can inject uninformative signal and distract the retriever.
Subsequently, adding the Cascade Reranking module recovers these losses (LAN: 84.55\% $\to$ 88.55\%; NO: 87.01\% $\to$ 90.38\%) and further increases overall accuracy to 88.42\% (a gain of 1.09 percentage points), suggesting that the reranker effectively filters the noisy candidates introduced by the naive Alpha fusion on text-dominant queries.
Finally, the Full R3-RAG, which integrates all modules including the reflection mechanism, reaches 91.51\% accuracy in this setting.
This final +3.09 percentage-point gain indicates that the self-reflection mechanism contributes to resolving retrieval failure scenarios and complex multi-step reasoning tasks.
The cumulative performance gain of 5.56 percentage points, together with the per-module pattern of degradation-then-recovery on text-centric subsets, suggests that the three modules provide complementary benefits: the Alpha module boosts visual recall but may introduce noise on text-only queries, which the reranker and reflection loop then correct.

\begin{table*}[t]
\centering
\renewcommand{\arraystretch}{1.2} 
\setlength{\abovecaptionskip}{0pt} 
\setlength{\belowcaptionskip}{10pt}

\caption{
\textbf{Ablation Study on ScienceQA.} 
We report the accuracy breakdown by Subject, Context Modality, and Grade. 
\textbf{Baseline} represents the standard RAG model.
The first three columns show the component settings.
}
\label{tab:ablation_full}


\vspace{0.01mm} 

\resizebox{1.0\textwidth}{!}{
\begin{tabular}{ccc ccc ccc cc c}
\toprule
\multicolumn{3}{c}{\textbf{Proposed Modules}} & \multicolumn{3}{c}{\textbf{Subject}} & \multicolumn{3}{c}{\textbf{Context Modality}} & \multicolumn{2}{c}{\textbf{Grade}} & \multirow{2}{*}{\textbf{Avg}} \\
\cmidrule(lr){1-3} \cmidrule(lr){4-6} \cmidrule(lr){7-9} \cmidrule(lr){10-11}
\textbf{Adapt. $\alpha$} & \textbf{Rerank} & \textbf{Reflect} & \textbf{NAT} & \textbf{SOC} & \textbf{LAN} & \textbf{TXT} & \textbf{IMG} & \textbf{NO} & \textbf{G1--6} & \textbf{G7--12} & \\
\midrule
% Row 1: Baseline
\textcolor{gray}{\ding{55}} & \textcolor{gray}{\ding{55}} & \textcolor{gray}{\ding{55}} & 86.50 & 82.23 & 87.82 & 86.31 & 82.90 & 90.03 & 88.81 & 80.81 & 85.95 \\
% Row 2: +Alpha
\ding{51} & \textcolor{gray}{\ding{55}} & \textcolor{gray}{\ding{55}} & 87.45 & 90.91 & 84.55 & 89.53 & 86.67 & 87.01 & 89.20 & 84.05 & 87.33 \\
% Row 3: +Alpha +Rerank
\ding{51} & \ding{51} & \textcolor{gray}{\ding{55}} & 87.30 & 91.11 & 88.55 & 88.59 & 86.96 & 90.38 & 91.26 & 83.32 & 88.42 \\
% Row 4: Full
\ding{51} & \ding{51} & \ding{51} & \textbf{91.34} & \textbf{92.35} & \textbf{91.18} & \textbf{89.86} & \textbf{90.98} & \textbf{93.17} & \textbf{93.25} & \textbf{88.40} & \textbf{91.51} \\
\bottomrule
\end{tabular}
}
\end{table*}
\section{Conclusion}\label{sec:conclusion}

This work presents R3-RAG, an adaptive framework that extends multimodal RAG with reweighting, reranking, and reflection modules for evidence-calibrated scientific reasoning.
Experimental findings suggest that the framework's effectiveness is related to its ability to manage retrieval noise and reduce the rigidity of one-shot retrieval pipelines.
Specifically, modality imbalance is mitigated through an intent-guided fusion mechanism, while cross-modal inconsistencies are reduced by an uncertainty-aware cascade reranker that uses score dispersion as a confidence signal.
Furthermore, by orchestrating a Diagnose-Rewrite-Route Reflection Loop, the system can iteratively revise low-confidence retrieval results.
Driven by these mechanisms, the framework achieves a competitive accuracy of 91.51\% on ScienceQA and shows preliminary generalization results on the MMMU test set official Science discipline subset.
Overall, this research provides empirical evidence that closed-loop retrieval refinement can improve scientific multimodal RAG under the evaluated settings.

At the same time, the present study has limitations that suggest future directions.
Because R3-RAG is evaluated mainly under a closed-domain retrieval setting, its effectiveness depends on the coverage and factual granularity of the local corpus; when specialized closed-book knowledge is absent from the corpus, retrieval may introduce distracting evidence rather than improve reasoning.
In addition, the current pipeline still incurs non-trivial inference costs, especially from final self-consistency generation, motivating more efficient uncertainty-triggered voting and selective reflection.
Broader evaluation across larger corpora, more diverse scientific domains, open-domain multimodal settings, and different backbone models is also needed to assess the scalability and generality of closed-loop evidence calibration.

% % To print the credit authorship contribution details
% \printcredits

%% Loading bibliography style file
\bibliographystyle{elsarticle-num}
%\bibliographystyle{cas-model2-names}

% Loading bibliography database
%\nocite{*}
\bibliography{cas-refs}



% Biography
%\bio{}
% Here goes the biography details.
%\endbio

%\bio{pic1}
% Here goes the biography details.
%\endbio

\end{document}
